\subsection{Density-one integrality without a birational lifting hypothesis}

\begin{claim}{D.25}{Proved}\label{clm:D25}
Let $U=\mathbf P^1_{\mathbf Q}-\{0,1,\infty\}$ and let
$s:G_{\mathbf Q}\to\pi_1(U)$ be a continuous section of the full
arithmetic fundamental group.  If $s$ is noncuspidal, its local
points are integral on
$\mathbf P^1_{\mathbf Z}-\{0,1,\infty\}$ at a set of primes of
Dirichlet density one.  They are also asymptotically integral:
for each positive integer $n$, all but finitely many ordinary
degenerate local points have boundary intersection multiplicity
divisible by $n$.  Equivalently, these positive multiplicities
tend to zero in $\widehat{\mathbf Z}$.

This conclusion does not assume that $s$ is birationally liftable.
It does not assert that the exceptional set is finite.
The same two conclusions hold for a noncuspidal continuous section
of any geometrically connected finite étale cover $X\to U$, with
respect to a good marked model of its smooth completion and boundary.
\end{claim}
\begin{proof}
We first prove the arithmetic fact needed to exclude the
exceptional alternative in the Legendre dichotomy:
\[
 \left.
 \begin{gathered}
 A\in\widehat{\mathbf Q^\times},\\
 A\bmod n=0\text{ in }\mathbf Q_p^\times/
       \mathbf Q_p^{\times n}
       \text{ for almost every }p,\text{ for each }n
 \end{gathered}
 \right\}
 \quad\Longrightarrow\quad A=1.
 \tag{D.25.1}
\]
The finite exceptional set is allowed to depend on $n$.
Represent $A\bmod n$ by $b_n\in\mathbf Q^\times$ and put
$F_n=\mathbf Q(\mu_n,b_n^{1/n})$.  Any nonidentity element of
$\operatorname{Gal}(F_n/\mathbf Q(\mu_n))$ acts on the roots of
$X^n-b_n$ by a nontrivial translation and fixes no root.
Chebotarev would then give infinitely many unramified primes at
which that polynomial has no local root.  This contradicts the
hypothesis.  Hence $b_n^{1/n}\in\mathbf Q(\mu_n)$.

After multiplication by a $2n$th root of unity, the positive real
radical $r=|b_n|^{1/n}$ lies in $\mathbf Q^{\mathrm{ab}}$.
Its field $\mathbf Q(r)$ is Galois over $\mathbf Q$ and lies in
$\mathbf R$.  Every conjugate of $r$ is therefore a real root of
$X^n-|b_n|$, which has at most two real roots.  Consequently
$r^2\in\mathbf Q$, and
$b_n^2=(r^2)^n\in\mathbf Q^{\times n}$.  It follows that $A^2=1$.
The valuation embedding
\[
 \widehat{\mathbf Q^\times}\hookrightarrow
 \{\pm1\}\times\prod_q\widehat{\mathbf Z}
\]
shows that its only torsion is $\{\pm1\}$.  The class $-1$
fails the hypothesis for $n=2$ at the infinitely many primes
$p\equiv3\pmod4$.  This proves (D.25.1).

Compatibility in $n$ is essential in this argument.  A single
level need not be detected outside a finite set: $16$ is an
eighth power in every $\mathbf Q_p$ with $p$ odd, but is not
an eighth power in $\mathbf Q$.  Here is a direct check.  For
$p\equiv3,5,7\pmod8$, the subgroup of eighth powers in
$\mathbf F_p^\times$ is respectively the subgroup of squares,
fourth powers, or squares, and contains $16$.  For
$p\equiv1\pmod8$, quadratic reciprocity gives a square root
$r$ of $2$ modulo $p$, and $r^8=16$.  Every root is nonzero,
so Hensel's lemma applies to the derivative $8r^7$ for odd
$p$.  A rational eighth root would have $2$-adic valuation
$1/2$, which is impossible.

By D.6 and the real section theorem
\cite[Corollary 3.13]{D_Mochizuki2003Real}, $s$ is a Selmer
section.  Apply Stix's Theorem 10 to the Legendre family
\[
 E_t:\quad y^2=x(x-1)(x-t)\longrightarrow U.
\]
The family has bad and not potentially good reduction at all
three boundary points.  The theorem asserts asymptotic
integrality and gives exactly two alternatives: either the
integral places have density one, or all but finitely many
places are cuspidal or degenerate and, for every $\ell$, the
representation pulled back by $s$ has the form
\[
 \rho_{s,\ell}\sim
 \begin{pmatrix}\delta\chi_\ell&*\\0&\delta\end{pmatrix},
 \tag{D.25.2}
\]
where $\delta$ is one quadratic character, independent of
$\ell$, ramified only at the finite set of bad places
\cite[Theorem 10]{D_Stix2015}.  We prove that the second
alternative makes $s$ cuspidal.

For every sufficiently good odd prime $p$, let
$c_p\in\{0,1,\infty\}$ be the cusp at which the local point
is cuspidal or into which it degenerates.  Work with $\ell=2$.
The local specialization calculation (equation (5.1) in the
proof of Stix's Theorem 20) says that $\rho_{s,2}(G_{\mathbf Q_p})$
cyclotomically normalizes the maximal unipotent group containing
the image of the geometric inertia at $c_p$, using its square
at $\infty$.  This is a calculation about the given local
point or cusp and the geometric pro-$2$ specialization; it
does not use a birational lift.  In the notation of that proof
we take $\beta=\operatorname{id}_U$ and remove no further
points.  The full arithmetic Legendre representation is already
congruent to the identity modulo $2$, because its full
$2$-torsion is rational over $U$, so the required pro-$2$
factorization is immediate.

It is important to distinguish the two local cases.  At an
ordinary degenerate point, after the quadratic splitting twist,
the inertia cocycle is $m_p$ times the tame $2$-adic character,
with $m_p>0$ an ordinary integer.  Its transvection is nonzero
and determines the Tate cyclotomic line.  At an exact cusp the
tangent cocycle may instead be a unit or zero, and the section's
arithmetic inertia image can be trivial.  The geometric cusp
inertia nevertheless has nontrivial unipotent image.  Since
the local section lies in its decomposition group, it
normalizes that geometric inertia through $\chi_2$.  Thus the
same cyclotomic-line assertion holds in the locally split
cuspidal case, without assuming nontrivial unipotent inertia
in the section image.

In (D.25.2) the unique maximal unipotent group normalized
through $\chi_2$ is the group fixing the global
$\delta\chi_2$ line.  The factor $\delta$ cancels in the
conjugation character.  If the local representation splits,
the other eigenspace gives $\chi_2^{-1}$, which differs from
$\chi_2$ because $\chi_2(\operatorname{Frob}_p)=p$ and
$p^2\neq1$ in $\mathbf Z_2$.  In the locally nonsplit case
the invariant line is unique.  This identifies the relevant
line independently of $p$.

The explicit geometric Legendre monodromies are
\[
 T_0=\begin{pmatrix}1&2\\0&1\end{pmatrix},\qquad
 T_1=\begin{pmatrix}1&0\\-2&1\end{pmatrix},\qquad
 T_\infty=\begin{pmatrix}1&-2\\2&-3\end{pmatrix}.
\]
Their fixed lines, with $T_\infty^2$ at infinity, reduce to
$(1:0),(0:1),(1:1)$ in $\mathbf P^1(\mathbf F_2)$.
Geometric conjugacy preserves those distinct reductions,
because the monodromy is trivial modulo $2$
\cite[Lemmas 17--19]{D_Stix2015}.  Hence $c_p$ is one fixed
cusp $c$ for all but finitely many $p$.  Only this first part
of Stix's proof of Theorem 20 is being used.  Its later
Riemann--Roch step, which needs lifts to additional open
subsets to distinguish arbitrary boundary points, is
unnecessary for these three cusps.

Choose the following global unit on $U$:
\[
 f=\begin{cases}
  1-t,&c=0,\\
  t,&c=1,\\
  t/(t-1),&c=\infty.
 \end{cases}
\]
It is regular at $c$ and has value $1$ there.  Put
$A=s^*\kappa(f)\in\widehat{\mathbf Q^\times}$.
At an ordinary point approaching $c$ at a sufficiently good
prime $p$, one has $f(x_p)\in1+p\mathbf Z_p$.  At a cusp
section at $c$, the Kummer pullback is the residue value $1$:
the function is regular and nonzero there, so its inertia
character is zero and its pullback is independent of the
tangent cocycle.  For each fixed $n$, excluding $p\mid n$
and using Hensel's lemma now shows that $A\bmod n$ is locally
trivial for all but finitely many $p$.  By (D.25.1), $A=1$.

Full local Kummer injectivity implies $f(x_p)=1$ at every
interior local point.  None of the displayed functions takes
the value $1$ on $U$; its unique such point on the completion
is $c$.  Thus no finite-local section is interior, and every
finite-local point of the completion lies in
$\{0,1,\infty\}$.  The finite-support argument in Stix's
Corollary 5 therefore makes $s$ cuspidal.  More precisely,
pass to finite étale neighbourhoods, use the finite inverse
image of this boundary as the support, and apply the corrected
finite-place form of Stoll's Theorem 8.2 already audited in
D.22.  That theorem only needs the prescribed finite support
at a density-one set of finite places; no assertion that the
real component is a boundary point is required.  The
neighbourhood limit gives a rational cusp pro-point.

The second alternative has therefore been excluded for a
noncuspidal $s$ on $U$.

Now let $h:X\to U$ be a geometrically connected finite étale
cover and let $s_X$ be noncuspidal.  Its pushforward $s_U$ is
noncuspidal as well.  Indeed, if $s_U$ lay in a cusp decomposition
group $D_c$, the intersection of $D_c$ with the corresponding
open subgroup $\pi_1(X)$ would be the cusp decomposition group
of a point $y$ of the normalized completion above $c$.
The section image projects onto all of $G_{\mathbf Q}$, so
the residue field of this $y$ must be $\mathbf Q$; the original
section would be cuspidal there.  The same assertion follows
after the geometric conjugacy used to specify $D_c$, which
merely chooses the lift of the cusp in the finite cover.

Normalize $\mathbf P^1_{\mathbf Z[1/S]}$ in $\mathbf Q(X)$
and enlarge the finite set $S$ so that the normalization and
its marked boundary give a good model, with the map finite
and étale over the model of $U$.  At any $p\notin S$ where
the local point of $s_U$ is integral in $U$, the point supplied
by D.6 for $s_X$ cannot be a cusp, because a cusp maps to a
cusp.  It is an interior $\mathbf Q_p$-point above the base
point and extends to the finite model over $\mathbf Z_p$ by
the valuative criterion for a finite morphism.  Thus it is
integral.  The set of such base primes has density one, and
discarding $S$ does not change that density.

Finally, at an ordinary degenerate place outside another
finite set, local boundary equations give
$d_p(U)=e_y d_p(X)$, where $e_y$ is the fixed horizontal
ramification index at the corresponding boundary point.
There are only finitely many possible $e_y$.  Divisibility
of $d_p(U)$ by $n\prod_y e_y$ at almost all such primes
therefore implies divisibility of $d_p(X)$ by $n$.
This proves asymptotic integrality upstairs, too.

Exact matrix products,
the three mod-$2$ lines and representative residue checks for
the single-level exception above are retained in
\path{code/D_legendre_density_certificate.py} and
\path{audit/D_legendre_density_certificate.json}.
\end{proof}

\begin{claim}{D.26}{Proved; number-field strengthening of D.25}\label{clm:D26}
Let $k$ be any number field. If
$A\in\widehat{k^\times}=\varprojlim_n k^\times/k^{\times n}$
is an $n$th power in $k_v$ at all but finitely many finite
places $v$, for every $n\ge1$, then $A=1$.
The exceptional finite set may depend on $n$.

Consequently, every noncuspidal continuous section of
$\mathbf P^1_k-\{0,1,\infty\}$, and of every geometrically
connected finite \'etale cover of this curve, has local points
integral at a set of finite places of Dirichlet density one.
They are asymptotically integral in the precise sense of D.25.
No birational lift or cofinite integrality is asserted.
\end{claim}
\begin{proof}
\emph{Cyclotomic reduction of a finite Kummer component.}
Fix a prime power $n=\ell^r$, represent the component of $A$
by $b\in k^\times$, and put
$F=k(\mu_n,b^{1/n})$.
If $\operatorname{Gal}(F/k(\mu_n))$ had a nonidentity element
$\tau$, it would act on the roots of $X^n-b$ by a nontrivial
translation and fix none. Its conjugates have the same
fixed-point property. Chebotarev gives infinitely many finite
places, unramified in $F$ and away from $n$ and the divisor of
$b$, with this Frobenius conjugacy class. At such a place
$X^n-b$ has no root in $k_v$, because a root in $k_v$ would
be fixed by its unramified local Frobenius. This contradicts
the hypothesis. Thus a chosen root $\alpha$ lies in $k(\mu_n)$.

Write $\Gamma=\operatorname{Gal}(k(\mu_n)/k)$. The Kummer
class is the inflation of
\[
 c(\sigma)=\frac{\sigma(\alpha)}{\alpha},\qquad
 c\in H^1(\Gamma,\mu_n).
 \tag{D.26.1}
\]
For each $\sigma\in\Gamma$, the same Chebotarev argument
shows that $\sigma$ fixes some root $\alpha\zeta$ of
$X^n-b$. Otherwise its Frobenius class would again give
infinitely many local failures. Hence
$c(\sigma)=\zeta/\sigma(\zeta)$, and the restriction of $c$
to the cyclic subgroup generated by $\sigma$ is a coboundary.
This conclusion holds for every cyclic subgroup of $\Gamma$.

\emph{The finite-level exceptional kernel has exponent at most two.}
If $\ell$ is odd, the group $\Gamma$ is a subgroup of the
cyclic group $(\mathbf Z/\ell^r\mathbf Z)^\times$.
Thus $\Gamma$ itself is cyclic, $c=0$, and
$b\in k^{\times\ell^r}$.
For $\ell=2$ and $r\ge3$, put
\[
 \Gamma_0=\Gamma\cap
       (1+4\mathbf Z/2^r\mathbf Z).
\]
The group $1+4\mathbf Z/2^r\mathbf Z$ is cyclic, and
$[\Gamma:\Gamma_0]\le2$. The restriction of $c$ to
$\Gamma_0$ vanishes. Inflation--restriction therefore puts
$c$ in the image of
\[
 H^1(\Gamma/\Gamma_0,\mu_{2^r}^{\Gamma_0})
       \longrightarrow H^1(\Gamma,\mu_{2^r}).
\]
Positive-degree cohomology of a finite group is killed by
its order, so $2c=0$. Equivalently,
$b^2\in k^{\times 2^r}$. For $r=1,2$ the group $\Gamma$
is cyclic and this exponent bound also holds.
Every odd-primary component of $A$ is consequently trivial,
and every $2$-primary component has square one. Chinese
remaindering in the power completion gives $A^2=1$.

\emph{Compatibility removes the exceptional torsion.}
The group $k^\times/\mu(k)$ is free abelian. Indeed, the
unit subgroup modulo roots of unity is finitely generated
free abelian by the unit theorem, the image of the divisor
map is a subgroup of the free abelian group on the finite
primes, and the resulting extension splits because its
quotient is free abelian. The power completion of a free
abelian group is torsion-free: its coordinate maps for a
basis embed it in a product of copies of
$\widehat{\mathbf Z}$. It follows that the torsion subgroup
of $\widehat{k^\times}$ is exactly $\mu(k)$. Its subgroup
killed by two is $\{1,-1\}$.

Choose $r$ so large that $\mu_{2^{r+1}}$ is not contained
in $k$. The roots of $X^{2^r}+1$ are exactly the primitive
$2^{r+1}$st roots of unity. Every nonidentity element of
$\operatorname{Gal}(k(\mu_{2^{r+1}})/k)$ fixes none of them:
it acts by a unit $u\ne1$ modulo $2^{r+1}$, and multiplying
$u-1$ by an odd exponent cannot give zero modulo $2^{r+1}$.
Chebotarev yields infinitely many finite places where $-1$
is not a $2^r$th power. Thus $A=-1$ is incompatible with
the hypothesis, and $A=1$.
This proves the first assertion without a single-level
local-to-global power assertion or a Grunwald--Wang input.

\emph{The density conclusion over $k$.}
We check explicitly the number-field substitutions in D.25.
Every finite completion $k_v$ is a finite extension of some
$\mathbf Q_p$, so D.6 and D.19 give the local point or cusp.
At real places the real section theorem gives the required
connected component, and complex places give no further
condition. Thus the section is Selmer in the sense used by
Stix. His Theorem 10 is stated over every number field and
applies to the same Legendre family
\cite[Theorem 10]{D_Stix2015}.

The full Legendre $2$-torsion is still rational over the
three-punctured line. The three geometric cusp monodromies
and their distinct reductions in
$\mathbf P^1(\mathbf F_2)$ are therefore exactly those of
D.25. In the triangular alternative, at a sufficiently
good finite place $v$ of odd residue characteristic,
\[
 \chi_2(\operatorname{Frob}_v)=N(v)>1,
\]
so $\chi_2$ and $\chi_2^{-1}$ are distinct on $G_{k_v}$.
The same unique cyclotomically normalized line determines
one of the three $k$-rational cusps at every such place,
and that cusp is independent of $v$. At an exact cusp,
the argument uses the nontrivial geometric inertia and its
normalization character, including when the tangent cocycle
makes the section's arithmetic inertia trivial. This is
the same locally split case checked in D.25.

Let $c$ be the resulting cusp and choose respectively
$f=1-t$, $t$, or $t/(t-1)$, so that $f$ is a unit on the
open curve, regular at $c$, with $f(c)=1$.
The class $s^*\kappa(f)$ lies in $\widehat{k^\times}$.
At an ordinary point degenerating into $c$ it is a principal
unit; at an exact cuspidal section its value is $1$,
independently of the tangent class. For each $n$, Hensel's
lemma away from the places dividing $n$ and the finite
model exception set makes its local class trivial at
almost every finite place. The first assertion just proved
gives $s^*\kappa(f)=1$.

Local Kummer injectivity over each $k_v$ then excludes every
finite-local interior point, since none of the three displayed
functions takes the value $1$ on the open curve. Stix's
Corollary 5, with the corrected finite-place Stoll argument
already checked in D.22, applies over every number field.
The finite support is the three-point boundary; no assertion
about a real point lying in that boundary is made. Hence the
triangular alternative is cuspidal. A noncuspidal section
must have density-one integral places, and asymptotic
integrality is the other conclusion of Theorem 10.

Finally let $X\to\mathbf P^1_k-\{0,1,\infty\}$ be a
geometrically connected finite \'etale cover. Intersecting a
cusp decomposition group with its open fundamental subgroup
shows, exactly as in D.25, that its noncuspidal sections
push forward to noncuspidal sections. Normalize
$\mathbf P^1_{\mathcal O_k[1/S]}$ in $k(X)$ and enlarge the
finite set $S$ to obtain a good marked model and a finite
map, \'etale above the model of the three-punctured line.
The finite valuative criterion lifts integral local points
integrally. Removing finitely many places does not change
Dirichlet density. At ordinary degenerate places the local
boundary multiplicities satisfy
$d_v(\text{base})=e_y d_v(X)$, with only finitely many
possible positive integers $e_y$. Divisibility downstairs
by $n\prod_y e_y$ implies divisibility upstairs by $n$.
This proves both assertions for the finite \'etale cover.
\end{proof}
